Research on navigation without, or with degraded, satellite positioning needs long recordings from low-cost inertial sensors, because the errors that matter grow over minutes to hours rather than seconds. Most public navigation datasets are instead short, collected on a campus or in a city, and built around cameras or lidar. MEMS-Nav contains \numRawRecordings{} drives recorded with the Sensor Logger application on \numDevices{} smartphone models (Google Pixel, Samsung Galaxy and Apple iPhone) between \numRawFirst{} and \numRawLast{} in the north-eastern United States, mostly at highway speed. The raw release keeps every sensor stream at the rate the phone delivered it, typically about 100~Hz for the accelerometer, gyroscope, magnetometer, operating-system gravity and attitude, and about 1~Hz for GNSS. The processed release cuts the recordings at inertial dropouts into \numTrajectories{} synchronized trajectories at 10~Hz, totalling \numTotalHours~h and \numTotalKm~km (median \numMedianHours~h and \numMedianKm~km), with \numFixes{} GNSS fixes left at their native rate and never interpolated. Each trajectory ships with cropped relief, free-air gravity anomaly and magnetic anomaly grids. The data are described per device, per channel and per trajectory, together with three properties of phone sensors that matter to anyone attempting map-based aiding: the receiver's advertised accuracies, a gravity vector whose norm the operating system holds at a device constant, and a magnetometer dominated by the vehicle's own field. The open-source simulator strapdown-rs replays every trajectory and can withhold or corrupt the GNSS fixes, and reference baselines are reported for an extended Kalman filter, an unscented Kalman filter and a Rao--Blackwellized particle filter. With every fix, the two Kalman filters agree with the recorded GNSS track to a median horizontal RMSE of \numEkfFullMedian~m and \numUkfFullMedian~m. With one corrupted fix per minute, the median grows to \numEkfDegradedMedian~m and \numUkfDegradedMedian~m. The reference position is the phone's own GNSS solution, not an independent ground truth.
